% BEGIN REV09_Q_COMPLEJO
\subsection{Itération du lecteur de rapports et récupération de phase}
\label{corpus9:iv:q}

Le lecteur \(R(x,y)=x/y^2\), appliqué aux représentations positives
de clôture et d’auto-échelle, admet la composition orientée
\[
 \mathcal Q_{\mathrm{rat}}(x,y)
   =\bigl(R(x,y),R(y,x)\bigr)
   =(x/y^2,y/x^2),\qquad x,y>0.
\]
Cette composition reçoit les coordonnées déjà produites par
APP--TRIT--TPK. Son étude conserve aussi bien le produit \(P=xy\)
que le rapport orienté \(O=x/y\), dont les transformations sont
\[
 P\circ\mathcal Q_{\mathrm{rat}}=P^{-1},\qquad
 O\circ\mathcal Q_{\mathrm{rat}}=O^3.
\]
La notation \(\mathcal Q_{\mathrm{rat}}\) distingue cette application
de l’extension aux intervalles du lecteur scalaire.

\begin{proposition}[Itération positive et prolongement complexe]
\label{corpus9:iv:q-iteracion}
Pour \(n\geq0\), soient
\[
 a_n=\frac{(-1)^n+3^n}{2},\qquad
 b_n=\frac{(-1)^n-3^n}{2}.
\]
Alors
\begin{equation}
 \mathcal Q_{\mathrm{rat}}^n(x,y)
   =(x^{a_n}y^{b_n},x^{b_n}y^{a_n}).
 \label{corpus9:iv:q-potencias}
\end{equation}
La même expression de Laurent définit une application sur
\((\mathbb C^\times)^2\). Pour \(n\geq1\), cette application est un
revêtement non ramifié de degré \(3^n\), de noyau
\begin{equation}
 \ker\mathcal Q_{\mathrm{rat}}^n
   =\{(\zeta,\zeta^{-1}):\zeta^{3^n}=1\}.
 \label{corpus9:iv:q-nucleo}
\end{equation}
Sa restriction au quadrant positif a une préimage positive
unique pour chaque paire positive.
\end{proposition}
\begin{proof}
Dans les coordonnées logarithmiques positives, la matrice est

\[
 B=\begin{pmatrix}1&-2\\-2&1\end{pmatrix}.
\]
Les vecteurs \((1,1)\) et \((1,-1)\) ont pour valeurs propres
\(-1\) et \(3\). La décomposition suivant ces deux directions donne
\(B^n=\left(
\begin{smallmatrix}a_n&b_n\\b_n&a_n\end{smallmatrix}\right)\)
et prouve \eqref{corpus9:iv:q-potencias}. Comme \(a_n,b_n\) sont
entiers, l’identité se prolonge comme identité de Laurent.


Écrivons \(m=3^n\), \(\epsilon=(-1)^n\). Pour une image
\((u,v)\in(\mathbb C^\times)^2\), le produit d’un antécédent
est fixé par \(p=xy=(uv)^\epsilon\). Les équations équivalent à
\[
 x^m=u\,p^{-b_n},\qquad y=p/x.
\]
La première a exactement \(m\) racines distinctes. La seconde
détermine \(y\) ; en outre, le produit des deux images est
\(p^\epsilon=uv\), de sorte que la seconde équation d’origine est
satisfaite. Toute fibre contient exactement \(m\) points. Pour
\((u,v)=(1,1)\), il vient \eqref{corpus9:iv:q-nucleo} ; les autres
fibres sont ses translatés dans le groupe multiplicatif.

La différentielle satisfait

\[
 \det D\mathcal Q_{\mathrm{rat}}^n
   =(-3)^n(xy)^{\epsilon-1}\ne0.
\]
La description locale par les racines d’un nombre non nul fournit
les \(m\) branches locales et prouve l’affirmation de revêtement.
Dans le quadrant positif, une seule de ces racines est positive.
En particulier, pour un pas,

\[
 x=(uv^2)^{-1/3},\qquad y=(u^2v)^{-1/3},
\]
avec les racines réelles positives.
\end{proof}

Dans le prolongement complexe, la condition \(xy=(uv)^{-1}\)
couple les deux racines cubiques de l’inverseur à un pas. Les
coordonnées \((P,O)\) identifient en outre
\((x,y)\) et \((-x,-y)\) ; leur application a pour noyau
\(\{(1,1),(-1,-1)\}\). La preuve précédente conserve \((x,y)\)
et, avec lui, la distinction entre cette identification de coordonnées
et les fibres de \(\mathcal Q_{\mathrm{rat}}^n\).

\begin{proposition}[Reconstruction par le registre ternaire
]
\label{corpus9:iv:q-acarreo}
Sur le sous-groupe

\[
 (x,y)=(e^{2\pi i\theta},e^{-2\pi i\theta}),
\]
l’application induit \(\theta\mapsto3\theta\pmod1\).
Les représentants \(\theta_j\in[0,1)\) étant fixés, soient
\[
 d_j=\lfloor3\theta_j\rfloor,\qquad
 \theta_{j+1}=3\theta_j-d_j,\qquad
 C_n=\sum_{j=0}^{n-1}d_j3^{n-1-j}.
\]
La récupération exacte est

\begin{equation}
 \theta_0=\frac{\theta_n+C_n}{3^n}.
 \label{corpus9:iv:q-recuperacion}
\end{equation}
\end{proposition}
\begin{proof}
La substitution dans l’expression de Laurent donne le triplement de
phase. La récurrence \(C_{n+1}=3C_n+d_n\), avec \(C_0=0\),
donne par récurrence \(\theta_n=3^n\theta_0-C_n\). Cela prouve
\eqref{corpus9:iv:q-recuperacion}. Chaque \(d_j\) appartient à
\(\{0,1,2\}\), de sorte que \(n\) chiffres ternaires conservent
le choix entre les \(3^n\) préimages de la phase réduite.
Les branches \(B_d(\theta)=(\theta+d)/3\) satisfont

\[
 B_a\circ B_b(\theta)=\frac{\theta+b+3a}{9}.
\]
Par conséquent, deux pas conservent le chiffre ordonné \(3d_0+d_1\).
\end{proof}

Le produit revient après deux itérations, tandis que le rapport
orienté est élevé à la neuvième puissance. La récupération démontrée
a pour domaine ce lecteur et son registre de phase. Son insertion
dans une représentation de l’état enrichi complet conserve
en outre les applications de feuillets, d’orientation, de mémoire et de
survie propres à cette représentation.
% END REV09_Q_COMPLEJO
